\chapter{Архитектура общего носителя четырёх слотов}
\label{ch:v9-four-slot-common-carrier}

\section{Endpoint-меню как внутреннее разложение}

Минимальный ранее допущенный endpoint-carrier имеет вид
\begin{equation}
 H_{24}=H_{21}\oplus H_{
m n}\oplus H_{
m t},
 \qquad \dim H_{
m n}=2,
 \qquad \dim H_{
m t}=1.
 \label{eq:v9-common-carrier-endpoint-splitting}
\end{equation}
Обозначим соответствующие ортогональные проекторы через
$P_{21},P_{
m n},P_{
m t}$. Тогда
\begin{equation}
 P_iP_j=\delta_{ij}P_i,
 \qquad P_{21}+P_{
m n}+P_{
m t}=I_{24},
 \qquad (\rank P_{21},\rank P_{
m n},\rank P_{
m t})=(21,2,1).
 \label{eq:v9-common-carrier-endpoint-projectors}
\end{equation}
Тем самым old, neutral-linking и family-triplet endpoint-секторы являются
подпространствами одного носителя, а не тремя сменяемыми Hilbert spaces.

\section{Почему требуется удвоенная цепная ячейка}

Шумовая ячейка Тома VIII уже имеет разложение
\begin{equation}
 K_{45}=K_0\oplus N_{42}\oplus E_{
m ch},
 \qquad (\dim K_0,\dim N_{42},\dim E_{
m ch})=(1,42,2).
 \label{eq:v9-common-carrier-noise-cell}
\end{equation}
Односторонний перенос имеет индекс $45$, тогда как локально допустимая
index-balanced ветвь требует компенсирующего обратного потока. Поэтому
общая локальная ячейка программы IX определяется как
\begin{equation}
 \Xi_{IX}=H_{24}\otimes K_{45}^{\rightarrow}
 \otimes K_{45}^{\leftarrow},
 \qquad \dim\Xi_{IX}=24\cdot45^2=48600.
 \label{eq:v9-common-carrier-xi}
\end{equation}
На одной и той же двусторонней цепи действуют две transport-ветви
\begin{equation}
 T_{
m F}=S_{45}^{\rightarrow}\otimes I,
 \qquad
 T_{\rm bal}=S_{45}^{\rightarrow}\otimes
 (S_{45}^{\leftarrow})^{-1}.
 \label{eq:v9-common-carrier-transport-menu}
\end{equation}
Их точные индексы равны
\begin{equation}
 \operatorname{ind}_{\rm GNVW}(T_{
m F})=45,
 \qquad
 \operatorname{ind}_{\rm GNVW}(T_{
m bal})=45\cdot45^{-1}=1.
 \label{eq:v9-common-carrier-transport-indices}
\end{equation}
Вторая ячейка является spectator в первой ветви и компенсатором во второй.

\section{Общая алгебра локальных данных}

Все локальные объекты теперь принадлежат одной алгебре
\begin{equation}
 \mathcal A_{IX}^{\rm cell}=\operatorname{End}(\Xi_{IX}).
 \label{eq:v9-common-carrier-algebra}
\end{equation}
Старый процесс вкладывается по карте
\begin{equation}
 X\longmapsto P_{21}XP_{21}\otimes|0\rangle\langle0|
 \otimes|0\rangle\langle0|,
 \label{eq:v9-common-carrier-old-embedding}
\end{equation}
а полный collision действует на первом environment-факторе:
\begin{equation}
 G(E_*,\chi)=\chi E_*G_0\otimes I_{45}^{\leftarrow},
 \qquad
 \mathcal L_{44}=\mathcal L_{42}
 +(\chi E_*)^2(\mathcal D_{L_-}+\mathcal D_{L_+}).
 \label{eq:v9-common-carrier-collision-family}
\end{equation}
Здесь формула задаёт общую operator-family, но не выбирает $E_*$ или
$\chi$.

Локальный положительный reference-parent можно записать без смены
пространства:
\begin{equation}
 h_{\Xi}=h_{24}\otimes I\otimes I
 +I\otimes h_{45}\otimes I
 +I\otimes I\otimes h_{45}\geq0.
 \label{eq:v9-common-carrier-reference-parent}
\end{equation}
Его форма совместима с обеими transport-ветвями, но его коэффициенты и
endpoint-ground sector ещё не выведены.

\section{Архитектурный реестр}

Машинный аудит проверяет восемь условий:
\begin{equation}
 N_{\rm common\ carrier\ architecture}=8/8.
 \label{eq:v9-common-carrier-architecture-ledger}
\end{equation}
При этом физический selector отсутствует:
\begin{equation}
 N_{\rm selected\ slots}=0/4,
 \qquad N_{\rm bounded\ four\ slot\ action}=0/1.
 \label{eq:v9-common-carrier-open-ledger}
\end{equation}

\begin{theorem}[Общий carrier четырёхслотовой программы]
Удвоенная цепь с локальной ячейкой
$\Xi_{IX}=H_{24}\otimes K_{45}^{\rightarrow}\otimes K_{45}^{\leftarrow}$
одновременно содержит все допущенные endpoint-секторы, полный
$44$-канальный collision и две transport-ветви с индексами $45$ и $1$.
Все объекты представлены в одной локальной алгебре, а старый процесс
восстанавливается точным угловым ограничением. Это закрывает архитектуру
общего носителя, но не выбирает четыре физических слота.
\end{theorem}

\begin{nogocriterion}[Общий носитель не является общим действием]
Прямая сумма endpoint-секторов не выбирает физический угол, а наличие двух
transport-операторов не выбирает закон переноса. Параметрическая семья
$G(E_*,\chi)$ также не является стационарным уравнением для $E_*$ и $\chi$.
Поэтому результат `8/8` запрещено читать как построение four-slot parent.
\end{nogocriterion}

\begin{protocol}\begin{enumerate}
\item endpoint carrier: \textbf{$H_{24}=21+2+1$};
\item noise cell: \textbf{$K_{45}=1+42+2$};
\item общий local carrier: \textbf{$\dim\Xi_{IX}=48600$};
\item old sector является углом: \textbf{да};
\item полный collision живёт на $\Xi_{IX}$: \textbf{да};
\item Floquet index: \textbf{$45$};
\item balanced index: \textbf{$1$};
\item архитектурный реестр: \textbf{$8/8$};
\item выбранные slots: \textbf{$0/4$};
\item bounded four-slot action: \textbf{$0/1$};
\item следующий гейт: \textbf{архитектура общего four-slot functional}.
\end{enumerate}\end{protocol}

Машинный сертификат сохранён в
\path{s2t_v9_four_slot_common_carrier_architecture_gate_results.json}.